\documentclass[11pt]{article}
\usepackage[margin=24mm]{geometry}
\usepackage{amsmath,amssymb,amsthm}
\usepackage[T1]{fontenc}
\usepackage{lmodern}
\usepackage[colorlinks=true,urlcolor=blue,linkcolor=blue,hypertexnames=false]{hyperref}
\newcommand{\HT}[2]{H_T^{#1}(#2)}
\newcommand{\Nm}{\operatorname{Nm}}
\newcommand{\Res}{\operatorname{Res}}
\newcommand{\Cor}{\operatorname{Cor}}
\newcommand{\Ind}{\operatorname{Ind}}
\newtheorem*{target}{Theorem II.3.11 (Tate's theorem)}
\newtheorem*{vanishing}{Theorem II.3.10 (support)}
\title{Tate's theorem for a class formation\\The splitting-module argument}
\author{J. S. Milne\\\small Source: \emph{Class Field Theory}, v4.03, 6 August 2020}
\date{}
\begin{document}
\maketitle
\noindent\small
Private reading transcription of the author's Theorem II.3.11, with the complete
supplied proof of its supporting Theorem II.3.10 and the separately published
correction to that proof. The original work is retained at
\url{https://www.jmilne.org/math/CourseNotes/CFT.pdf}.
The selected material occupies physical PDF pages 92--94 (printed pages 83--85).
The correction is retained with the official errata page in the source package.
Mathematical expressions and the author's argument are retained; connective prose
and layout are normalized. This is published-material transcription, not
author-native TeX or a new proof.
\normalsize

% BEGIN CONTEXT TATE CONVENTIONS
\section*{Notation and the author's convention}
The group $G$ is finite. In this part of Section II.3 the author uses Tate
cohomology throughout, dropping the subscript $T$ except in the main statements.
We preserve that convention in the transcribed proofs. In positive degrees these
groups agree with ordinary group cohomology. In degree zero the convention is
\[
 \HT{0}{H,M}=M^H/\Nm_HM,\qquad
 \Nm_H(m)=\sum_{h\in H}hm.
\]
The action on $\mathbb Z$ is trivial. The notation $M_0$ in the supporting
argument denotes the underlying abelian group with trivial action, as in the
source's preceding induced-module conventions (II.1.3--1.4).
The preceding cohomological tools used by the author are cyclic periodicity,
the exact cohomology sequence, the double boundary map (II.1.13),
restriction/corestriction (II.1.33), inflation/restriction (II.1.34),
and the Tate-acyclicity of induced modules. These are context for the supplied
construction, rather than additional counted problems.
% END CONTEXT TATE CONVENTIONS

% BEGIN PRIMARY STATEMENT II.3.11
\section*{The target statement}
\begin{target}
Let $G$ be a finite group and $C$ a $G$-module. Suppose that for every subgroup
$H\leq G$, including $H=G$,
\begin{enumerate}
\item $H^1(H,C)=0$;
\item $H^2(H,C)$ is cyclic of order $(H:1)$.
\end{enumerate}
Then for every integer $r$ there is an isomorphism
\[
 \HT{r}{G,\mathbb Z}\longrightarrow\HT{r+2}{G,C}
\]
depending only on the choice of a generator for $H^2(G,C)$.
\end{target}
% END PRIMARY STATEMENT II.3.11

% BEGIN SUPPORT II.3.10 ORIGINAL
\section*{The supporting all-degree vanishing theorem}
\begin{vanishing}
Let $G$ be a finite group and $M$ a $G$-module. If
\[
 H^1(H,M)=0=H^2(H,M)
\]
for all subgroups $H$ of $G$, then $\HT{r}{G,M}=0$ for all $r\in\mathbb Z$.
\end{vanishing}
\begin{proof}[Author proof as printed on p.83]
For cyclic $G$ the conclusion follows from periodicity.
For solvable $G$, argue by induction on its order. There is a proper normal
subgroup $H$ for which $G/H$ is cyclic. The pair $(H,M)$ satisfies the
hypotheses, so induction gives $H^r(H,M)=0$ for all $r$. By II.1.34 there are
exact sequences, for $r\geq1$,
\[
 0\longrightarrow H^r(G/H,M^H)\longrightarrow H^r(G,M)
   \longrightarrow H^r(H,M). \tag{*}
\]
Because $H^1(G,M)=H^2(G,M)=0$, these give
$H^1(G/H,M^H)=H^2(G/H,M^H)=0$. Cyclic periodicity for $G/H$ then gives
$H^r(G/H,M^H)=0$ for all $r$. The sequences $(*)$ imply
$H^r(G,M)=0$ for $r>0$.

To show $H^0(G,M)=0$, let $x\in M^G$. Since $H^0(G/H,M^H)=0$, there is
$y\in M^H$ such that $\Nm_{G/H}(y)=x$. Since $H^0(H,M)=0$, there is
$z\in M$ with $\Nm_H(z)=y$. Hence
\[
 \Nm_G(z)=(\Nm_{G/H}\circ\Nm_H)(z)=x.
\]
Thus $H^r(G,M)=0$ for every $r\geq0$.

To proceed further, tensor the augmentation sequence (19) with $M$, obtaining
\[
 0\longrightarrow M'\longrightarrow\mathbb Z[G]\otimes_{\mathbb Z}M
  \longrightarrow M\longrightarrow0.
\]
By II.1.4 and II.1.3, respectively,
\[
 \mathbb Z[G]\otimes_{\mathbb Z}M
 \simeq\mathbb Z[G]\otimes_{\mathbb Z}M_0
 \simeq\Ind^G(M_0)
\]
is induced. Therefore $H^r(H,\mathbb Z[G]\otimes_{\mathbb Z}M)=0$ for every
$r$ and every subgroup $H$ of $G$, and the cohomology sequence gives
\[
 H^r(H,M)\simeq H^{r+1}(H,M')
\]
for all $r$ and all $H$. In particular, $M'$ satisfies the hypotheses of the
theorem. The part already proved gives $H^r(G,M')=0$ for $r\geq0$, so
\[
 0=H^0(G,M')=H^{-1}(G,M).
\]
Repeating this argument gives $H^{-2}(G,M)=0$, and so on. This proves the
theorem for solvable $G$.

For an arbitrary finite $G$, its Sylow $p$-subgroup $G_p$ and $M$ also satisfy
the hypotheses. The printed proof concludes that $H^r(G_p,M)=0$ for all $r$
and $p$, and, using II.1.33, that the $p$-primary component of $H^r(G,M)$
is zero for all $r$ and $p$. It then concludes $H^r(G,M)=0$ for all $r$.
\end{proof}
% END SUPPORT II.3.10 ORIGINAL

% BEGIN SUPPORT II.3.10 AUTHOR ERRATUM
\subsection*{Separately published correction to II.3.10}
The official errata page records Kiran Kedlaya's correction to the preceding
proof. The use of $H^r$ in place of $H_T^r$ obscures the degree-zero issue in
the reduction to solvable groups. That reduction as invoked in the printed
Sylow argument is valid for $r>0$, while degree zero is also needed for
shifting to negative degrees.

The correction supplies the missing degree-zero step.
Since $\HT{0}{G_p,M}=0$, the group $\HT{0}{G,M}$ is killed by $[G:G_p]$.
Varying $p$ gives $\HT{0}{G,M}=0$; dimension shifting then gives the
negative-degree conclusions. This paragraph records the correction supplied
on the author's errata page; it is separate from the printed proof above.
% END SUPPORT II.3.10 AUTHOR ERRATUM

% BEGIN PRIMARY PROOF II.3.11 SPLITTING MODULE
\section*{The author's proof of II.3.11}
\begin{proof}[Author proof]
Choose a generator $\gamma$ of $H^2(G,C)$. Because
$\Cor\circ\Res=(G:H)$, the element $\Res(\gamma)$ generates $H^2(H,C)$
for every subgroup $H$ of $G$.

Let $\varphi$ be a cocycle representing $\gamma$. Define $C(\varphi)$ as the
direct sum of $C$ and the free abelian group with basis $x_\sigma$, one for
each $\sigma\in G$, $\sigma\ne1$. Extend the $G$-action on $C$ by setting
\[
 \sigma x_\tau=x_{\sigma\tau}-x_\sigma+\varphi(\sigma,\tau).
\]
Here $x_1$ is interpreted as $\varphi(1,1)$, exactly as in the source.
This defines an action, since
\[
 (\rho\sigma)x_\tau=x_{\rho\sigma\tau}-x_{\rho\sigma}
                  +\varphi(\rho\sigma,\tau),
\]
whereas
\begin{align*}
 \rho(\sigma x_\tau)
 &=\rho\bigl(x_{\sigma\tau}-x_\sigma+\varphi(\sigma,\tau)\bigr)\\
 &=x_{\rho\sigma\tau}-x_\rho+\varphi(\rho,\sigma\tau)
  -\bigl(x_{\rho\sigma}-x_\rho+\varphi(\rho,\sigma)\bigr)
  +\rho\varphi(\sigma,\tau).
\end{align*}
These expressions agree by the cocycle condition
\[
 \rho\varphi(\sigma,\tau)+\varphi(\rho,\sigma\tau)
 =\varphi(\rho\sigma,\tau)+\varphi(\rho,\sigma).
\]
The cocycle $\varphi$ is the coboundary of the $1$-cochain
$\sigma\mapsto x_\sigma$, so $\gamma$ maps to zero in $H^2(G,C(\varphi))$.
For this reason $C(\varphi)$ is called the splitting module for $\gamma$.
% END PRIMARY PROOF II.3.11 SPLITTING MODULE

% BEGIN PRIMARY PROOF II.3.11 SUBGROUP VANISHING
We first show that the hypotheses imply
\[
 H^1(H,C(\varphi))=0=H^2(H,C(\varphi))
\]
for every subgroup $H$ of $G$. The augmentation sequence is
\[
 0\longrightarrow I_G\longrightarrow\mathbb Z[G]
  \longrightarrow\mathbb Z\longrightarrow0,
\]
where $I_G$ is free abelian on $\sigma-1$, with $\sigma\in G$, $\sigma\ne1$.
The module $\mathbb Z[G]$ is induced, so $H^r(H,\mathbb Z[G])=0$ for all $r$.
Consequently
\begin{align*}
 H^1(H,I_G)&\simeq H^0(H,\mathbb Z)\simeq\mathbb Z/(H:1)\mathbb Z,\\
 H^2(H,I_G)&\simeq H^1(H,\mathbb Z)=0.
\end{align*}
Define the additive map $\alpha:C(\varphi)\longrightarrow\mathbb Z[G]$ by
\[
 \alpha(c)=0\quad(c\in C),\qquad\alpha(x_\sigma)=\sigma-1.
\]
The resulting sequence of $G$-modules
\[
 0\longrightarrow C\longrightarrow C(\varphi)
   \xrightarrow{\alpha}I_G\longrightarrow0
\]
is exact. Its cohomology sequence reads
\[
\begin{split}
0\longrightarrow H^1(H,C(\varphi))\longrightarrow H^1(H,I_G)
 \longrightarrow H^2(H,C)\\
 \xrightarrow{\,0\,}H^2(H,C(\varphi))\longrightarrow0.
\end{split}
\]
The end zeros use $H^1(H,C)=0$ and $H^2(H,I_G)=0$.
The indicated map from $H^2(H,C)$ is zero because that group is generated
by $\Res(\gamma)$, whose image is the restriction of the image of $\gamma$
in $H^2(G,C(\varphi))$, already zero.
Thus $H^1(H,I_G)\longrightarrow H^2(H,C)$ is onto, and is an isomorphism
because both groups have the same finite order. Its kernel and cokernel,
$H^1(H,C(\varphi))$ and $H^2(H,C(\varphi))$, both vanish.
% END PRIMARY PROOF II.3.11 SUBGROUP VANISHING

% BEGIN PRIMARY PROOF II.3.11 DOUBLE BOUNDARY
By Theorem II.3.10, $H^r(H,C(\varphi))=0$ for all $r$.
Splicing the two short exact sequences gives
\[
 0\longrightarrow C\longrightarrow C(\varphi)\longrightarrow\mathbb Z[G]
   \longrightarrow\mathbb Z\longrightarrow0. \tag{24}
\]
Both middle terms are Tate-acyclic: for every $r$,
$H^r(G,C(\varphi))=0=H^r(G,\mathbb Z[G])$.
The double boundary map is therefore an isomorphism, by II.1.13,
\[
 H^r(G,\mathbb Z)\longrightarrow H^{r+2}(G,C).
\]
\end{proof}
% END PRIMARY PROOF II.3.11 DOUBLE BOUNDARY

\medskip
\noindent\small
The single collected target is II.3.11. The supporting theorem, its published
correction, and the auxiliary exact sequences add no separate problem count.
The generator is chosen, rather than asserted canonical. The source's following
Remark II.3.13 identifies the map as cup product with that chosen element;
this reader retains the supplied splitting-module proof and does not add an
independent proof of the remark or the later local-field application.
\end{document}
